\frametitle{Schr\"odinger picture}
\begin{small}
{\scriptsize
The time-dependent Schr\"odinger equation (or equation of motion) reads
\[
\imath \hbar\frac{\partial }{\partial t}|\Psi_S(t)\rangle = \hat{H}\Psi_S(t)\rangle,
\]
where the subscript $S$ stands for Schr\"odinger here.
A formal solution is given by 
\[
|\Psi_S(t)\rangle = \exp{(-\imath\hat{H}(t-t_0)/\hbar)}|\Psi_S(t_0)\rangle.
\]
The Hamiltonian $\hat{H}$ is hermitian and the exponent represents a unitary 
operator with an operation carried ut on the wave function at a time $t_0$.
}
\end{small}
}
\frame
{
\frametitle{Interaction picture}
\begin{small}
{\scriptsize
Our Hamiltonian is normally written out as the sum of an unperturbed part $\hat{H}_0$ and an interaction part $\hat{H}_I$, that is
\[
\hat{H}=\hat{H}_0+\hat{H}_I.
\]
In general we have $[\hat{H}_0,\hat{H}_I]\ne 0$ since $[\hat{T},\hat{V}]\ne 0$.
We wish now to define a unitary transformation in terms of $\hat{H}_0$ by defining
\[
|\Psi_I(t)\rangle = \exp{(\imath\hat{H}_0t/\hbar)}|\Psi_S(t)\rangle,
\]
which is again a unitary transformation carried out now at the time $t$ on the 
wave function in the Schr\"odinger picture. 
}
\end{small}
}
\frame
{
\frametitle{Interaction picture}
\begin{small}
{\scriptsize
We can easily find the equation of motion by taking the time derivative
\[
\imath \hbar\frac{\partial }{\partial t}|\Psi_I(t)\rangle = -\hat{H}_0\exp{(\imath\hat{H}_0t/\hbar)}\Psi_S(t)\rangle+\exp{(\imath\hat{H}_0t/\hbar)}
\imath \hbar\frac{\partial }{\partial t}\Psi_S(t)\rangle.
\]
}
\end{small}
}
\frame
{
\frametitle{Interaction picture}
\begin{small}
{\scriptsize
Using the definition of the Schr\"odinger equation, we can rewrite the last equation as 
\[
\imath \hbar\frac{\partial }{\partial t}|\Psi_I(t)\rangle = \exp{(\imath\hat{H}_0t/\hbar)}\left[-\hat{H}_0+\hat{H}_0+\hat{H}_I\right]\exp{(-\imath\hat{H}_0t/\hbar)}\Psi_I(t)\rangle,
\]
which gives us
\[
\imath \hbar\frac{\partial }{\partial t}|\Psi_I(t)\rangle = \hat{H}_I(t)\Psi_I(t)\rangle,
\]
 with 
\[
\hat{H}_I(t)=
\exp{(\imath\hat{H}_0t/\hbar)}\hat{H}_I\exp{(-\imath\hat{H}_0t/\hbar)}.
\]
}
\end{small}
}
\frame
{
\frametitle{Interaction picture}
\begin{small}
{\scriptsize
The order of the operators is important since $\hat{H}_0$ and $\hat{H}_I$ do generally not commute.
The expectation value of
an arbitrary operator in the interaction picture can now be written as
\[
\langle \Psi'_S(t)|\hat{O}_S|\Psi_S(t)\rangle = 
\langle \Psi'_I(t) |\exp{(\imath\hat{H}_0t/\hbar)}\hat{O}_I
\exp{(-\imath\hat{H}_0t/\hbar)}|\Psi_I(t)\rangle,
\]
and using the definition
\[
\hat{O}_I(t)=
\exp{(\imath\hat{H}_0t/\hbar)}\hat{O}_I\exp{(-\imath\hat{H}_0t/\hbar)},
\]
we obtain
\[
\langle \Psi'_S(t)|\hat{O}_S|\Psi_S(t)\rangle = 
\langle \Psi'_I(t) |\hat{O}_I(t)|\Psi_I(t)\rangle,
\]
stating that a unitary transformation does not change expectation values!
}
\end{small}
}
\frame
{
\frametitle{Interaction picture}
\begin{small}
{\scriptsize
If the take the time derivative of the operator in the interaction picture we arrive at the following equation of motion
\[
\imath \hbar\frac{\partial }{\partial t}\hat{O}_I(t) = \exp{(\imath\hat{H}_0t/\hbar)}\left[\hat{O}_S\hat{H}_0-\hat{H}_0\hat{O}_S\right]\exp{(-\imath\hat{H}_0t/\hbar)}=\left[\hat{O}_I(t),\hat{H}_0\right].
\]
Here we have used the time-independence of the Schr\"odinger equation
together with the observation that any function of an operator commutes with the operator itself. 
}
\end{small}
}
\frame
{
\frametitle{Interaction picture}
\begin{small}
{\scriptsize
In order to solve the equation of motion equation in the interaction picture, we define a unitary operator
time-development operator $\hat{U}(t,t')$. Later we will derive its
connection with the linked-diagram theorem, which yields a
linked expression for the actual operator. 
The action of the operator on the wave function is
\[
|\Psi_I(t) \rangle = \hat{U}(t,t_0)|\Psi_I(t_0)\rangle,
\]
with the obvious value $\hat{U}(t_0,t_0)=1$.
}
\end{small}
}
\frame
{
\frametitle{Interaction picture}
\begin{small}
{\scriptsize
The time-development operator $U$ has the
properties that
\[
     \hat{U}^{\dagger}(t,t')\hat{U}(t,t')=\hat{U}(t,t')\hat{U}^{\dagger}(t,t')=1,
\]
which implies that $U$ is unitary
\[
     \hat{U}^{\dagger}(t,t')=\hat{U}^{-1}(t,t').
\]
Further,
\[
    \hat{U}(t,t')\hat{U}(t't'')=\hat{U}(t,t'')
\]
and
\[
    \hat{U}(t,t')\hat{U}(t',t)=1,
\]
which leads to
\[
    \hat{U}(t,t')=\hat{U}^{\dagger}(t',t).
\]
}
\end{small}
}
\frame
{
\frametitle{Interaction picture}
\begin{small}
{\scriptsize
Using our definition of Schr\"odinger's equation in the interaction picture, we can then construct the operator $\hat{U}$. We have defined
\[
|\Psi_I(t)\rangle = \exp{(\imath\hat{H}_0t/\hbar)}|\Psi_S(t)\rangle,
\]
which can be rewritten as 
\[
|\Psi_I(t)\rangle = \exp{(\imath\hat{H}_0t/\hbar)}\exp{(-\imath\hat{H}(t-t_0)/\hbar)}|\Psi_S(t_0)\rangle,
\]
or
\[
|\Psi_I(t)\rangle = \exp{(\imath\hat{H}_0t/\hbar)}\exp{(-\imath\hat{H}(t-t_0)/\hbar)}\exp{(-\imath\hat{H}_0t_0/\hbar)}|\Psi_I(t_0)\rangle.
\]
}
\end{small}
}
\frame
{
\frametitle{Interaction picture}
\begin{small}
{\scriptsize
From the last expression we can define
\[
\hat{U}(t,t_0)=\exp{(\imath\hat{H}_0t/\hbar)}\exp{(-\imath\hat{H}(t-t_0)/\hbar)}\exp{(-\imath\hat{H}_0t_0/\hbar)}.
\]
It is then easy to convince oneself that the properties defined above are satisfied by the definition of $\hat{U}$. 
}
\end{small}
}
\frame
{
\frametitle{Interaction picture}
\begin{small}
{\scriptsize
We derive the equation of motion for $\hat{U}$ using the above definition.
This results in
\[
\imath \hbar\frac{\partial }{\partial t}\hat{U}(t,t_0) = \hat{H}_I(t)\hat{U}(t,t_0),
\]
which we integrate from $t_0$ to a time $t$ resulting in
\[
\hat{U}(t,t_0)-\hat{U}(t_0,t_0)=\hat{U}(t,t_0)-1=-\frac{\imath}{\hbar}\int_{t_0}^t dt' \hat{H}_I(t')\hat{U}(t',t_0),
\]
which can be rewritten as
\[
\hat{U}(t,t_0)=1-\frac{\imath}{\hbar}\int_{t_0}^t dt' \hat{H}_I(t')\hat{U}(t',t_0).
\]
}
\end{small}
}
\frame
{
\frametitle{Interaction picture}
\begin{small}
{\scriptsize
We can solve this equation iteratively keeping in mind the time-ordering of the of the operators
\[
\hat{U}(t,t_0)=1-\frac{\imath}{\hbar}\int_{t_0}^t dt' \hat{H}_I(t')+\left(\frac{-\imath}{\hbar}\right)^2\int_{t_0}^t dt'\int_{t_0}^{t'} dt'' \hat{H}_I(t')\hat{H}_I(t'')+\dots
\]
The third term can be written as 
\[
\int_{t_0}^t dt'\int_{t_0}^{t'} dt'' \hat{H}_I(t')\hat{H}_I(t'')=
\frac{1}{2}\int_{t_0}^t dt'\int_{t_0}^{t'} dt'' \hat{H}_I(t')\hat{H}_I(t'')
+\frac{1}{2}\int_{t_0}^t dt''\int_{t''}^{t} dt' \hat{H}_I(t')\hat{H}_I(t'').
\]
}
\end{small}
}
\frame
{
\frametitle{Interaction picture}
\begin{small}
{\scriptsize
We obtain this expression by changing the integration order in the second term
via a change of the integration variables $t'$ and $t''$  in 
\[
\frac{1}{2}\int_{t_0}^t dt'\int_{t_0}^{t'} dt'' \hat{H}_I(t')\hat{H}_I(t'').
\]
We can rewrite the terms which contain the double integral as
\[
\int_{t_0}^t dt'\int_{t_0}^{t'} dt'' \hat{H}_I(t')\hat{H}_I(t'')=\]
\[
\frac{1}{2}\int_{t_0}^t dt'\int_{t_0}^{t'} dt''\left[\hat{H}_I(t')\hat{H}_I(t'')\Theta(t'-t'')
+\hat{H}_I(t')\hat{H}_I(t'')\Theta(t''-t')\right],
\]
with $\Theta(t''-t')$ being the standard Heavyside or step function. The step function allows us to give a specific time-ordering to the above expression.
}
\end{small}
}
\frame
{
\frametitle{Interaction picture}
\begin{small}
{\scriptsize
With the $\Theta$-function we can rewrite the last expression as 
\[
\int_{t_0}^t dt'\int_{t_0}^{t'} dt'' \hat{H}_I(t')\hat{H}_I(t'')=
\frac{1}{2}\int_{t_0}^t dt'\int_{t_0}^{t'} dt''\hat{T}\left[\hat{H}_I(t')\hat{H}_I(t'')\right],
\]
where $\Hat{T}$ is the so-called time-ordering operator. 
}
\end{small}
}
\frame
{
\frametitle{Interaction picture}
\begin{small}
{\scriptsize
With this definition, we can rewrite the expression for $\hat{U}$ as 
\[
\hat{U}(t,t_0)=\sum_{n=0}^{\infty}\left(\frac{-\imath}{\hbar}\right)^n\frac{1}{n!}
\int_{t_0}^t dt_1\dots \int_{t_0}^t dt_N \hat{T}\left[\hat{H}_I(t_1)\dots\hat{H}_I(t_n)\right]=\hat{T}\exp{\left[\frac{-\imath}{\hbar}
\int_{t_0}^t dt' \hat{H}_I(t')\right]}.
\]
The above time-evolution operator in the interaction picture will be used
to derive various contributions to many-body perturbation theory. See also exercise 26 for a discussion of the various time orderings.
}
\end{small}
}
\frame
{
\frametitle{Heisenberg picture}
\begin{small}
{\scriptsize
We wish now to define a unitary transformation in terms of $\hat{H}$ by defining
\[
|\Psi_H(t)\rangle = \exp{(\imath\hat{H}t/\hbar)}|\Psi_S(t)\rangle,
\]
which is again a unitary transformation carried out now at the time $t$ on the 
wave function in the Schr\"odinger picture. If we combine this equation with 
Schr\"odinger's equation we obtain the following equation of motion
\[
\imath \hbar\frac{\partial }{\partial t}|\Psi_H(t)\rangle = 0,
\]
meaning that $|\Psi_H(t)\rangle$ is time independent. An operator in this picture is defined as
\[
\hat{O}_H(t)=
\exp{(\imath\hat{H}t/\hbar)}\hat{O}_S\exp{(-\imath\hat{H}t/\hbar)}.
\]
}
\end{small}
}
\frame
{
\frametitle{Heisenberg picture}
\begin{small}
{\scriptsize
The time dependence is then in the operator itself, and this yields in turn the
following equation of motion
\[
\imath \hbar\frac{\partial }{\partial t}\hat{O}_H(t) = \exp{(\imath\hat{H}t/\hbar)}\left[\hat{O}_H\hat{H}-\hat{H}\hat{O}_H\right]\exp{(-\imath\hat{H}t/\hbar)}=\left[\hat{O}_H(t),\hat{H}\right].
\]
We note that an operator in the Heisenberg picture can be related to the corresponding
operator in the interaction picture as 
\[
\hat{O}_H(t)=
\exp{(\imath\hat{H}t/\hbar)}\hat{O}_S\exp{(-\imath\hat{H}t/\hbar)}=\]
\[
\exp{(\imath\hat{H}_It/\hbar)}\exp{(-\imath\hat{H}_0t/\hbar)}\hat{O}_I\exp{(\imath\hat{H}_0t/\hbar)}\exp{(-\imath\hat{H}_It/\hbar)}.
\]
}
\end{small}
}
\frame
{
\frametitle{Heisenberg picture}
\begin{small}
{\scriptsize
With our definition of the time evolution operator we see that
\[
\hat{O}_H(t)=\hat{U}(0,t)\hat{O}_I\hat{U}(t,0),
\]
which in turn implies that $\hat{O}_S=\hat{O}_I(0)=\hat{O}_H(0)$, all operators are equal at $t=0$. The wave function in the Heisenberg formalism is 
related to the other pictures as 
\[
|\Psi_H\rangle=|\Psi_S(0)\rangle=|\Psi_I(0)\rangle,
\]
since the wave function in the Heisenberg picture is time independent. 
We can relate this wave function to that a given time $t$ via the time evolution operator as
\[
|\Psi_H\rangle=\hat{U}(0,t)|\Psi_I(t)\rangle.
\]
}
\end{small}
}
\frame
{
\frametitle{Adiabatic hypothesis}
\begin{small}
{\scriptsize
We assume that the interaction term is switched on gradually. Our wave function at time $t=-\infty$ and $t=\infty$ is supposed to represent a non-interacting system
given by the solution to the unperturbed part of our Hamiltonian $\hat{H}_0$.
We assume the ground state is given by $|\Phi_0\rangle$, which could be a Slater determinant.
We define our Hamiltonian as
\[
\hat{H}=\hat{H}_0+\exp{(-\varepsilon t/\hbar)}\hat{H}_I,
\]
where $\varepsilon$ is a small number. The way we write the Hamiltonian 
and its interaction term is meant to simulate the switching of the interaction.
}
\end{small}
}
\frame
{
\frametitle{Adiabatic hypothesis}
\begin{small}
{\scriptsize
The time evolution of the wave function in the interaction picture is then
\[
|\Psi_I(t) \rangle = \hat{U}_{\varepsilon}(t,t_0)|\Psi_I(t_0)\rangle,
\]
with 
\[
\hat{U}_{\varepsilon}(t,t_0)=\sum_{n=0}^{\infty}\left(\frac{-\imath}{\hbar}\right)^n\frac{1}{n!}
\int_{t_0}^t dt_1\dots \int_{t_0}^t dt_N \exp{(-\varepsilon(t_1+\dots+t_n)/\hbar)}\hat{T}\left[\hat{H}_I(t_1)\dots\hat{H}_I(t_n)\right]
\]
}
\end{small}
}
\frame
{
\frametitle{Adiabatic hypothesis}
\begin{small}
{\scriptsize
In the limit $t_0\rightarrow -\infty$, the solution ot Schr\"odinger's equation is
$|\Phi_0\rangle$, and the eigenenergies are given by 
\[
\hat{H}_0|\Phi_0\rangle=W_0|\Phi_0\rangle,
\]
meaning that 
\[
|\Psi_S(t_0)\rangle = \exp{(-\imath W_0t_0/\hbar)}|\Phi_0\rangle,
\]
with the corresponding interaction picture wave function given by
\[
|\Psi_I(t_0)\rangle = \exp{(\imath \hat{H}_0t_0/\hbar)}|\Psi_S(t_0)\rangle=|\Phi_0\rangle.
\]
}
\end{small}
}
\frame
{
\frametitle{Adiabatic hypothesis}
\begin{small}
{\scriptsize
The solution becomes time independent in the limit $t_0\rightarrow -\infty$.
The same conclusion can be reached by looking at 
\[
\imath \hbar\frac{\partial }{\partial t}|\Psi_I(t)\rangle =
\exp{(\varepsilon |t|/\hbar)}\hat{H}_I|\Psi_I(t)\rangle 
\]
and taking the limit $t\rightarrow -\infty$.
We can rewrite the equation for the wave function at a time $t=0$ as
\[
|\Psi_I(0) \rangle = \hat{U}_{\varepsilon}(0,-\infty)|\Phi_0\rangle.
\]
}
\end{small}
}
